%%%%%%%%%%%%%%%%%%%%%%%%%%%%%%%%%%%%%%%%%%%%%%%
%
% Template per Elaborato di Laurea
% DISI - Dipartimento di Ingegneria e Scienza dell’Informazione
%
% update 2015-09-10
%
% Per la generazione corretta del 
% pdflatex nome_file.tex
% bibtex nome_file.aux
% pdflatex nome_file.tex
% pdflatex nome_file.tex
%
%%%%%%%%%%%%%%%%%%%%%%%%%%%%%%%%%%%%%%%%%%%%%%%

% to compile for the version to print, uncomment \printtrue
\newif\ifprint
  % \printtrue

\ifprint
  \documentclass[epsfig,a4paper,11pt,titlepage,twoside,openright]{book}
\else
  % formato FRONTE RETRO
  \documentclass[epsfig,a4paper,11pt,titlepage,twoside,openany]{book}
\fi
\usepackage{epsfig}
\usepackage{plain}
\usepackage{setspace}
\usepackage[paperheight=29.7cm,paperwidth=21cm,outer=1.5cm,inner=2.5cm,top=2cm,bottom=2cm]{geometry} % per definizione layout
\usepackage{titlesec} % per formato custom dei titoli dei capitoli
\usepackage{enumitem}
\usepackage{minted}
\usemintedstyle{pastie}
\usepackage{hyperref}
\usepackage{xcolor}
\definecolor{LightGray}{gray}{0.95}
\usepackage{amsmath}
\usepackage{placeins}
\usepackage{graphicx}
\graphicspath{ {./assets/} }

%% Calligraphic alphabet
\newcommand*{\cA}{\ensuremath{\mathcal{A}}}
\newcommand*{\cB}{\ensuremath{\mathcal{B}}}
\newcommand*{\cC}{\ensuremath{\mathcal{C}}}
\newcommand*{\cD}{\ensuremath{\mathcal{D}}}
\newcommand*{\cE}{\ensuremath{\mathcal{E}}}
\newcommand*{\cF}{\ensuremath{\mathcal{F}}}
\newcommand*{\cG}{\ensuremath{\mathcal{G}}}
\newcommand*{\cH}{\ensuremath{\mathcal{H}}}
\newcommand*{\cI}{\ensuremath{\mathcal{I}}}
\newcommand*{\cJ}{\ensuremath{\mathcal{J}}}
\newcommand*{\cK}{\ensuremath{\mathcal{K}}}
\newcommand*{\cL}{\ensuremath{\mathcal{L}}}
\newcommand*{\cM}{\ensuremath{\mathcal{M}}}
\newcommand*{\cN}{\ensuremath{\mathcal{N}}}
\newcommand*{\cO}{\ensuremath{\mathcal{O}}}
\newcommand*{\cP}{\ensuremath{\mathcal{P}}}
\newcommand*{\cQ}{\ensuremath{\mathcal{Q}}}
\newcommand*{\cR}{\ensuremath{\mathcal{R}}}
\newcommand*{\cS}{\ensuremath{\mathcal{S}}}
\newcommand*{\cT}{\ensuremath{\mathcal{T}}}
\newcommand*{\cU}{\ensuremath{\mathcal{U}}}
\newcommand*{\cV}{\ensuremath{\mathcal{V}}}
\newcommand*{\cW}{\ensuremath{\mathcal{W}}}
\newcommand*{\cX}{\ensuremath{\mathcal{X}}}
\newcommand*{\cY}{\ensuremath{\mathcal{Y}}}
\newcommand*{\cZ}{\ensuremath{\mathcal{Z}}}

%% program, formats,...
\newcommand*{\vmtlib}{\texttt{VMT-LIB}}
\newcommand*{\smtlib}{\texttt{SMT-LIB}}
\newcommand*{\vmt}{\texttt{VMT}}
\newcommand*{\btor}{\texttt{BTOR2}}
\newcommand*{\VMT}{Verification Modulo Theory}
\newcommand*{\smt}{\texttt{SMT}}
\newcommand*{\sat}{\texttt{SAT}}
\newcommand*{\SMT}{Satisfiability Modulo Theory}
\newcommand*{\smv}{\texttt{SMV}}
\newcommand*{\ltl}{\texttt{LTL}}
\newcommand*{\LTL}{Linear Temporal Logic}
\newcommand*{\boolector}{\texttt{Boolector}}
\newcommand*{\pysmt}{\texttt{pySMT}}
\newcommand*{\pyvmt}{\texttt{PyVmt}}
\newcommand*{\nuxmv}{\texttt{nuXmv}}
\newcommand*{\iceia}{\texttt{ic3ia}}
\newcommand*{\euforia}{\texttt{EUForia}}
\newcommand*{\nusmv}{\texttt{nuSMV}}
\newcommand*{\yices}{\texttt{Yices}}
\newcommand*{\zthree}{\texttt{Z3}}
\newcommand*{\mathsat}{\texttt{MathSAT}}
\newcommand*{\python}{Python}
\newcommand*{\pytest}{PyTest}

%% websites, organizations,...
\newcommand*{\fbk}{FBK}
\newcommand*{\FBK}{Fondazione Bruno Kessler}
\newcommand*{\github}{GitHub}
\newcommand*{\pypi}{PyPi}
\newcommand*{\ieee}{IEEE}

%% LTL
\newcommand*{\F}{\ensuremath{\mathit{F}}}
\newcommand*{\G}{\ensuremath{\mathit{G}}}
\newcommand*{\R}{\ensuremath{\mathit{R}}}
\newcommand*{\U}{\ensuremath{\mathit{U}}}
\newcommand*{\X}{\ensuremath{\mathit{X}}}

\newcommand*{\fF}{\ensuremath{\F{} \phi}}
\newcommand*{\fG}{\ensuremath{\G{} \phi}}
\newcommand*{\fR}{\ensuremath{\phi \R{} \psi}}
\newcommand*{\fU}{\ensuremath{\phi \U{} \psi}}
\newcommand*{\fX}{\ensuremath{\X{} \phi}}

\newcommand*{\fsafe}{\ensuremath{\G{} \phi}}
\newcommand*{\ffair}{\ensuremath{\G{} (\F{} \phi)}}
\newcommand*{\flive}{\ensuremath{\F{} (\G{} \phi)}}

%% styles
\newcommand{\inputmintedpy}[3]{\inputminted[
    firstline=#2, lastline=#3, frame=lines,
    framesep=2mm, baselinestretch=1.2, bgcolor=LightGray,
    fontsize=\footnotesize, linenos]{python3}{#1}
}


%%%%%%%%%%%%%%
% supporto lettere accentate
%
%\usepackage[latin1]{inputenc} % per Windows;
\usepackage[utf8x]{inputenc} % per Linux (richiede il pacchetto unicode);
%\usepackage[applemac]{inputenc} % per Mac.

\singlespacing

\usepackage[english]{babel}

\begin{document}

  % nessuna numerazione
  \pagenumbering{gobble} 
  \pagestyle{plain}

\thispagestyle{empty}

\begin{center}
  \begin{figure}[h!]
    \centerline{\psfig{file=marchio_unitrento_colore_it_202002.eps,width=0.6\textwidth}}
  \end{figure}

  \vspace{2 cm} 

  \LARGE{Department of Information Engineering and Computer Science\\}

  \vspace{1 cm} 
  \Large{Bachelor degree in\\
  Computer Science}

  \vspace{2 cm} 
  \Large\textsc{final dissertation\\} 
  \vspace{1 cm} 
  \Huge\textsc{\pyvmt{}\\}
  \Large{\it{A \python{} Library to manage \VMT{} Problems}}


  \vspace{2 cm} 
  \begin{tabular*}{\textwidth}{ c @{\extracolsep{\fill}} c }
  \Large{Supervisor} & \Large{Student}\\
  \Large{Sebastiani Roberto}& \Large{Girardi Alessandro}\\
  \\
  \Large{Co-supervisor} & \\
  \Large{Griggio Alberto}& \\
  \end{tabular*}

  \vspace{2 cm} 

  \Large{Academic year 2021/2022}
  
\end{center}



  \clearpage
 
%%%%%%%%%%%%%%%%%%%%%%%%%%%%%%%%%%%%%%%%%%%%%%%%%%%%%%%%%%%%%%%%%%%%%%%%%%
%%%%%%%%%%%%%%%%%%%%%%%%%%%%%%%%%%%%%%%%%%%%%%%%%%%%%%%%%%%%%%%%%%%%%%%%%%
%% Nota
%%%%%%%%%%%%%%%%%%%%%%%%%%%%%%%%%%%%%%%%%%%%%%%%%%%%%%%%%%%%%%%%%%%%%%%%%%
%% Sezione Ringraziamenti opzionale
%%%%%%%%%%%%%%%%%%%%%%%%%%%%%%%%%%%%%%%%%%%%%%%%%%%%%%%%%%%%%%%%%%%%%%%%%%
%%%%%%%%%%%%%%%%%%%%%%%%%%%%%%%%%%%%%%%%%%%%%%%%%%%%%%%%%%%%%%%%%%%%%%%%%%
  % \thispagestyle{empty}

\begin{center}
  {\bf \Huge Acknowledgements}
\end{center}

\vspace{4cm}


\emph{
  ...thanks to...
}

  % \clearpage
  \pagestyle{plain} % nessuna intestazione e pie pagina con numero al centro

  
  % inizio numerazione pagine in numeri arabi
  \mainmatter

%%%%%%%%%%%%%%%%%%%%%%%%%%%%%%%%%%%%%%%%%%%%%%%%%%%%%%%%%%%%%%%%%%%%%%%%%%
%%%%%%%%%%%%%%%%%%%%%%%%%%%%%%%%%%%%%%%%%%%%%%%%%%%%%%%%%%%%%%%%%%%%%%%%%%
%% Nota
%%%%%%%%%%%%%%%%%%%%%%%%%%%%%%%%%%%%%%%%%%%%%%%%%%%%%%%%%%%%%%%%%%%%%%%%%%
%% Si ricorda che il numero massimo di facciate e' 30.
%% Nel conteggio delle facciate sono incluse 
%%   indice
%%   sommario
%%   capitoli
%% Dal conteggio delle facciate sono escluse
%%   frontespizio
%%   ringraziamenti
%%   allegati    
%%%%%%%%%%%%%%%%%%%%%%%%%%%%%%%%%%%%%%%%%%%%%%%%%%%%%%%%%%%%%%%%%%%%%%%%%%
%%%%%%%%%%%%%%%%%%%%%%%%%%%%%%%%%%%%%%%%%%%%%%%%%%%%%%%%%%%%%%%%%%%%%%%%%%

    % indice
    \tableofcontents
    \clearpage
    
    
          
    % gruppo per definizone di successione capitoli senza interruzione di pagina
    \begingroup
      \ifprint
      \else
        % nessuna interruzione di pagina tra capitoli
        % ridefinizione dei comandi di clear page
        \renewcommand{\cleardoublepage}{} 
        \renewcommand{\clearpage}{} 
      \fi
      % redefinizione del formato del titolo del capitolo
      % da formato
      %   Capitolo X
      %   Titolo capitolo
      % a formato
      %   X   Titolo capitolo
      
      \titleformat{\chapter}
        {\normalfont\Huge\bfseries}{\thechapter}{1em}{}
        
      \titlespacing*{\chapter}{0pt}{0.59in}{0.02in}
      \titlespacing*{\section}{0pt}{0.20in}{0.02in}
      \titlespacing*{\subsection}{0pt}{0.10in}{0.02in}
      
      % sommario
      \chapter*{Abstract}
\label{abstract}

\addcontentsline{toc}{chapter}{Abstract}

Formal verification has become an essential tool in ensuring safety and liveness of complex software and hardware systems that are used in safety critical situations.
Over the last few years symbolic methods for formal verification have become quite popular, and the definition of the standard language \smtlib{} has led to the creation of a large collection of benchmarks.
This allowed for comparisons between \smt{} solvers which further sparked innovation.
Another contributor is \pysmt{}, a \python{} library that makes working with \smt{} problems simple, and offers tools for formula manipulation which don't rely on any specific solver.
This thesis goal is to contribute to the field of \VMT{} by extending \pysmt{} with the creation of \pyvmt{}, a library that includes specific tools used in verification, allowing for faster prototyping of models and algorithms, and increasing the interoperability of model checkers.
I will be going over the architecture of the library, the ways it allows to run model checkers within the library and how it interfaces with external tools.
I will also be presenting the results from an experimental evaluation, showing how the implementations of two \ltl{} encoding procedures compare to each other as a proof of concept for the capabilities of the library.


%%%%%%%%%%%%%%%%%%%%%%%%%%%%%%%%%%%%%%%%%%%%%%%%%%%%%%%%%%%%%%%%%%%%%%%%%%
%%%%%%%%%%%%%%%%%%%%%%%%%%%%%%%%%%%%%%%%%%%%%%%%%%%%%%%%%%%%%%%%%%%%%%%%%%
%% Nota
%%%%%%%%%%%%%%%%%%%%%%%%%%%%%%%%%%%%%%%%%%%%%%%%%%%%%%%%%%%%%%%%%%%%%%%%%%
%% Sommario e' un breve riassunto del lavoro svolto dove si descrive 
%% l’obiettivo, l’oggetto della tesi, le metodologie e 
%% le tecniche usate, i dati elaborati e la spiegazione delle conclusioni 
%% alle quali siete arrivati.
%% Il sommario dell’elaborato consiste al massimo di 3 pagine e deve contenere le seguenti informazioni: 
%%   contesto e motivazioni
%%   breve riassunto del problema affrontato
%%   tecniche utilizzate e/o sviluppate
%%   risultati raggiunti, sottolineando il contributo personale del laureando/a
%%%%%%%%%%%%%%%%%%%%%%%%%%%%%%%%%%%%%%%%%%%%%%%%%%%%%%%%%%%%%%%%%%%%%%%%%%
%%%%%%%%%%%%%%%%%%%%%%%%%%%%%%%%%%%%%%%%%%%%%%%%%%%%%%%%%%%%%%%%%%%%%%%%%%      
      
      %%%%%%%%%%%%%%%%%%%%%%%%%%%%%%%%
      % lista dei capitoli
      %
      % \input oppure \include
      %
      \chapter{Introduction}
\label{introduction}

The classical approach for finding and resolving bugs in software usually involves writing tests and code reviews done by other developers.
The safety of a system tested in this manner is not guaranteed since only a limited amount of scenarios can be explored, and subtle bugs may go unnoticed even after peer-review.
A mathematical approach is therefore necessary whenever such guarantees are required, for example when dealing with complex, safety-critical systems where bugs may have severe repercussions and are therefore not acceptable.

There are multiple approaches to mathematically formalizing a program, one of them is representing it as a transition system, that is, a model utilizing states and transitions between the states.
Transition systems can also be thought of as directed graphs, where each node is a state and the branches connecting the states are the possible transitions.
After representing a program in a formal manner, it is possible to specify some property which we want the system to respect and, because of the mathematical formalization, proving such specifications becomes possible using some specialized program called model checker.
% For these applications, software may be written utilizing a formal language, allowing for model checking to be applied in order to formally verify some properties corresponding to the specifications.

\paragraph*{Futurebus+}
The verification of the \ieee{} Futurebus+ cache coherence protocol \cite{DBLP:journals/fmsd/ClarkeGHJLMN95} is regarded as the first time formal methods were applied to verify the correctness of an \ieee{} protocol.
The approach to verification required constructing a model using a formal language (in this case \smv{} was chosen) and using it to show that the system satisfied the specifications.
This research was also able to find errors within the original formulation that were previously unnoticed, thus demonstrating a practical industrial usage of formal methods in the definition of standards.

\section{Contribution}

\pyvmt{} was created by me as part of a formative internship at \FBK{}. It's an extension of an existing project which was also created in the foundation, \pysmt{}, which offers a solver-agnostic formula manipulation interface.
Several factors motivate the creation for an extension:
\begin{itemize}
  \item the creation of specific tools and functions which are commonly used in model verification
  \item the extension of the existing \smtlib{} parser to extract the information relating to the model and its properties written in \vmtlib{} format
  \item the addition of support for model checkers to offer an experience similar to \pysmt{} while dealing with verification problems
\end{itemize}

The objective that \pyvmt{} poses for itself is creating a library to deal with verification problems that does not rely on any specific model checker, not just being the wrapper for one specific standard but adding tools and abstractions which further simplify interoperability between systems and sharing of benchmarks.

\paragraph*{What can \pyvmt{} do?}
\pyvmt{} is a \python{} library that enables the user to manage \vmt{} models by reading and writing code in the \vmtlib{} format, modifying models by adding variables, constraints, properties, or with some higher level function such as renaming and synchronous composition. It offers several model checkers as backend for verification, making it simpler for the user to choose several of them and compare their performance.

The features of the library will be presented in an order which aids comprehension of the architecture of the library, but first in Chapter \ref{ch:background} I will go over the scientific and technical background, explaining in detail what \pysmt{} is and the importance it has in its field.
In Chapter \ref{ch:modelling} the modelling system of \pyvmt{} will be explained in detail with examples showcasing the API for model manipulation.
Chapters \ref{ch:vmtlib} and \ref{ch:model-checking} will explain how this library can interface with external systems thanks to model parsing, serialization, and built-in functions which allow the user to run model checkers and get their results.
In Chapter \ref{ch:ltl-support} I will demonstrate some of the possibilities which are enabled by \pyvmt{} through an experimental evaluation which serves as a proof of concept for the capabilities of this library to implement complex algorithms.
Finally, in Chapter \ref{ch:other-features} other utility features will be presented and in Chapter \ref{ch:conclusions} I will conclude with some final remarks and lay the base for future works.

\paragraph*{Used techniques.}
\pyvmt{} is written in \python{}, it was developed using a continuous integration approach, utilizing \github{} actions as a tool to verify non-regression after each update.
Tests are run using the \pytest{} test suite and are mostly composed of unit tests which aim to cover most of the library's source code.
Most of the documentation is present in the source code as doc-strings so that they can easily be looked up while using an IDE. An external reference is also available and can be generated using sphinx.
After all the tests pass, the new version of the library is automatically uploaded as a package to \pypi{} so that any user can use it in their own project.

% \paragraph*{Achieved results.}
% To verify the functionality of \pyvmt{}, I collected several benchmarks and ran them against the available solvers. The results from these experiments show a promising support from the backends to deal with most of the problems within their supported theories.


      \chapter{Verification modulo theory}

\paragraph*{SMT}
Satisfiability modulo theory is a generalization of the SAT (boolean satisfiability) problem. It involves proving whether a formula can be satisfied by some assignment of the variables. The expressions, which can include lists, arrays, bit vectors, and strings are interpreted within a formal theory. There exists a standard language for defining SMT problems called SMT-LIB \cite{SMT-LIB} and it simplifies interfacing to the solvers, which are the programs developed to decide whether a formula is satisfiable.

\paragraph*{VMT}
Verification modulo theory adds a dynamic aspect to SMT, allowing the specification of transition systems. In VMT, transition systems are defined in a symbolic manner, they're composed of several elements:
\begin{itemize}
    \item State variables
    \item An init constraint
    \item A transition constraint
\end{itemize}
States are therefore not explicitly expressed, rather the assignments to the state variables may be used to encode them.
This has the advantage of allowing for transition systems with an extremely large number of states, and is frequently a more natural way to define their behavior.
The transition between states is regulated by a transition constraint, which is a function between current and next state variables and must be verified during a transition. The set of initial states is also not explicit, initial states must verify the init constraint.


\section{Importance of model checking}
The classical approach for finding and resolving bugs in software usually involves writing tests and code reviews done by other developers.
The safety of a system tested in this manner is not guaranteed since only a limited amount of scenarios can be explored, and subtle bugs may go unnoticed even after peer-review.
A mathematical approach is therefore necessary whenever such guarantees are required, for example when dealing with complex, safety-critical systems where bugs may have severe repercussions and are therefore not acceptable.

For these applications, software may be written utilizing a formal language, allowing for model checking to be applied in order to formally verify some properties corresponding to the specifications.

\section{Some applications}
Many different problems can be solved by using a verification tool, the following paragraphs include a few relevant ones.

\paragraph*{LTL} LTL or Linear Temporal Logic is used to define complex properties over execution sequences of the system.
It works with a fragment of first order logic and adds several operators (X, F, G, U) to introduce the concept of time to our properties.
The semantics of the operators is as follows:
\begin{itemize}
    \item \begin{math} X \phi \end{math}, in the next step \begin{math} \phi \end{math} is verified
    \item \begin{math} F \phi \end{math}, in some future step \begin{math} \phi \end{math} has to be verified
    \item \begin{math} G \phi \end{math}, \begin{math} \phi \end{math} is globally verified
    \item \begin{math} \phi\ R\ \psi \end{math}, \begin{math} \psi \end{math} is verified as long as \begin{math} \phi \end{math} is not verified
\end{itemize}

Combinations of these operators may be used to create classic property types or more complex properties.

\paragraph*{Safety} Safety properties are used to define that the system never reaches states in which a condition we want to avoid is verified, in LTL this can be written as \begin{math}G (\lnot bad)\end{math}.
This verifies that the states which the system can reach is limited to safe ones and that the bad behavior cannot be reached during the execution.

An example of a safety problem is ensuring that it's impossible for the door of a microwave to be open when the device is running:
\begin{math}
    G \lnot (open \land running)
\end{math}.

\paragraph*{Fairness} Fairness specifies that a set of states must be visited infinitely often during execution and is denoted in LTL as \begin{math} G(F(\phi))\end{math}.
Fairness can be used both as a property to be checked or as a constraint that defines how infinite traces of the system should look like.

A problem involving a fairness property is ensuring that a scheduler for an operating system eventually schedules each process, thus not leaving any process to wait infinitely often.

\paragraph*{Liveness} Liveness properties specify that after a while, a property holds forever. It's denoted in LTL as \begin{math} F(G(\phi)) \end{math}.
This kind of property is frequently explained with the notion that something good will eventually happen and relates to the concept of progress.
They're similar to fairness properties, in fact the previous formula is equivalent to \begin{math} \lnot G(F(\lnot \phi))\end{math} meaning that a liveness property may be proven invalid by a trace in which \begin{math} \phi \end{math} is negated infinitely often, which is a fair path for \begin{math}\lnot \phi\end{math}.

\paragraph*{Decision problems} Some complex logic problems involving decisions on an evolving system may be modelled as a transition system.
After that is done, a property can be added specifying that a solution to the original problem may not be found, i.e. \begin{math} G (\lnot solved) \end{math}.
A specialized program may be used to verify whether the property holds, if it doesn't (meaning that a solution to our problem exists), it may print the exact steps required to reach such solution.

A classic problem that can be solved in this manner is the ferryman problem, where a ferryman has to carry across a river several entities (e.g. a cabbage, a wolf, and a sheep) on a boat with a limited capacity, while never leaving certain pairs of entities unattended (e.g. the sheep with the cabbage or the sheep with the wolf).
Once modelled, the solution can be found as a counterexample to the property:
\begin{math}
    G (\exists (entity). \lnot isAcross(entity))
\end{math}.
The resulting counterexample displays which entities have to be carried across the river and back in each step.

\section{The VMT-LIB format}
VMT-LIB is an extension of the SMT2 format and provides a standard interface for defining transition systems by exploiting the possibility to attach annotations to terms and formulas, which can be used to specify the model behaviors. \cite{VMT-LIB}

It is designed with a limited amount of language constructs, aiming for simplicity for anyone interested in implementing the language within their VMT solver:
\begin{itemize}[noitemsep]
    \item \textbf{:next name} is used to specify state variables. To define a state variable, two variables are created, one representing the current state and one representing the next state. The next annotation is used to link the next state variable to the current state variable.
    \item \textbf{:init} is used to specify the formula used to constrain the initial states for the system.
    \item \textbf{:trans} is used to specify the formula used to constrain the transitions from each state to the next for the system.
    \item \textbf{:invar-property idx} is used to specify an invariant property to be verified, in LTL it's in the form \begin{math} G \phi \end{math} and requires a formula to always hold.  A unique non-negative integer specifies the index of the property.
    \item \textbf{:live-property idx} is used to specify a live property to be verified, in LTL it's in the form \begin{math} F\ G\ \phi\end{math} and requires a formula to hold infinitely often. A unique non-negative integer specifies the index of the property.
\end{itemize}
\cite{vmt-lib-paper} %...

% execution paths...


      \chapter{Modelling using PyVmt}
PyVmt exports a Model class which allows the user to define a model from scratch, manually adding variables and formulae or using helper functions to do so.

\section{State variables}
\section{Inputs}
\section{Constraints}
\section{Properties}

      \chapter{VMT-LIB support}

\section{The VMT-LIB format}
VMT-LIB is an extension of the SMT2 format and provides a standard interface for defining transition systems by exploiting the possibility to attach annotations to terms and formulas, which can be used to specify the model behaviors. \cite{vmt-lib-paper}

It is designed with a limited amount of language constructs, aiming for simplicity for anyone interested in implementing the language within their VMT solver:
\begin{itemize}
    \item \textbf{:next name} is used to specify state variables. To define a state variable, two variables are created, one representing the current state and one representing the next state. The next annotation is used to link the next state variable to the current state variable.
    \item \textbf{:init} is used to specify the formula used to constrain the initial states for the system.
    \item \textbf{:trans} is used to specify the formula used to constrain the transitions from each state to the next for the system.
    \item \textbf{:invar-property idx} is used to specify an invariant property to be verified, in LTL it's in the form \begin{math} G \phi \end{math} and requires a formula to always hold.  A unique non-negative integer specifies the index of the property.
    \item \textbf{:live-property idx} is used to specify a live property to be verified, in LTL it's in the form \begin{math} F\ G\ \phi\end{math} and requires a formula to hold infinitely often. A unique non-negative integer specifies the index of the property.
\end{itemize}

An in-depth explanation of the format, including a list of available tools may be found on the VMT-LIB website. \cite{VMT-LIB}

PyVmt modelling system is similar to that of VMT-LIB making converting to and from the format simpler.
This allows for higher interoperability since a model created in PyVmt may be serialized to VMT-LIB and be used in a model checker which isn't supported by the library.

\section{Reading}
PySMT already supports parsing SmtLib2 through the SmtLibParser class.
PyVmt defines a function for reading a script from an input stream called `read', which parses the script and runs through the annotations.
For every NEXT annotation a new state variable is added. Since PyVmt only keeps track of curr state variables the next state one is ignored.
After that, any declaration which was not a curr or next state variable is considered as an input.

To make next state variables work properly in trans constraints and properties, a substituter is employed to replace every instance of a next state variable with the corresponding representation using the Next operator.

Reading an existing VMT-LIB model can have several use-cases:
\begin{enumerate}
    \item Modifying a model by adding state variables or constraints in a programmatic way, or combining it with other models
    \item Analyzing an existing model (e.g. by creating a custom formula walker that reads the constraints, by printing some statistics,\dots)
    \item Running the same VMT-LIB script against multiple solvers, with different combination of options
    \item PySmt's SmtLibParser shows exactly where parsing errors occur, helping debug broken VMT-LIB files. This can be demoed by running \mintinline{bash}{python3 -m pyvmt < some_script.vmt}
\end{enumerate}
\inputminted[firstline=11, lastline=34]{python3}{py/vmtlib_support.py}

\section{Writing}
PySMT already provides printers for SMT-LIB2 supporting two formats: normal and daggified, the latter providing a version that avoids repeating repeated formulae.

The serialize method of the Model produces a fresh symbol for each of the current state variables so that it may use them to replace all of the Next operators which are not supported in VMT-LIB.
To make sure that the replacements work properly, a custom walker named NextPusher is used to push all of the Next operators down to the leaves, leaving Next operators only on Symbols.
Additionally, to replace the unsupported operators the VmtLibSubstituter is used. To be specific it performs a replacement of any Next state variable with the corresponding fresh symbol, leaving the final formula with only supported operators.

As previously mentioned, serializing a model to VMT-LIB is useful whenever trying to use a model checker which is not directly supported by PyVmt, but it can also be a step to use a model checker that doesn't support VMT-LIB in the first place.
Once the model has been serialized to VMT-LIB, several tools are available to further convert the model to a different format, they may be found on the vmt-lib website under the section `Tool Support' \cite{VMT-LIB}, the currently supported formats for conversion are:
\begin{itemize}
    \item btor2
    \item horn
    \item nuXmv
    \item core vmt, excluding language extensions
\end{itemize}

\inputminted[firstline=37, lastline=38]{python3}{py/vmtlib_support.py}

\section{LTL}
PyVmt supports reading LTL-properties thanks to an extension of the SmtLibParser that adds the various operators as interpreted functions. Additionally, LTL operators may be printed using the VmtPrinter and VmtDagPrinter, which respectively extend the SmtPrinter and SmtDagPrinter provided by pySMT. The \begin{math}\phi\ R\ \psi\end{math} operator is not natively supported by VMT-LIB, properties that contain it must first be passed through the VmtLibSubstituter which replaces it with the equivalent \begin{math}\lnot (\lnot \phi\ U \lnot \psi) \end{math}. As mentioned in the previous section this is done automatically while serializing a model.

% LTL removal


      \chapter{Model checking in \pyvmt{}}
\label{ch:model-checking}
\pyvmt{} supports several model checkers and provides a layer of abstraction to interface with them. Through a model checker it is possible to know whether the properties specified over a model are verified.

Running a model checker over a model in \pyvmt{} can be done entirely within \python{} code through the following steps:
\begin{enumerate}
    \item Read or create a new model
    \item Create a new instance of the preferred solver
    \item Optionally modify the options for the solver
    \item Formulate a property and either add it to the model or pass it directly to one of the helper methods available for the model checker
    \item Analyze the result
\end{enumerate}

\begin{listing}[h]
    \label{alg:model-checking-steps}
    \inputmintedpy{py/model_checking.py}{5}{20}
    \caption{A snippet of code showcasing the steps which can be used to read a model from the standard input and check its first property using \iceia{}.}
\end{listing}

Typically to run a solver, \pyvmt{} creates a temporary file and stores in it the model with the property it's trying to check.
After that, it creates a subprocess using \texttt{subprocess.run}, and passing the options either through command line arguments or through the standard input.
Finally, the output from the model checker is parsed and returned to the user.

\section{Available solvers}

\paragraph*{\nuxmv{}.}
\nuxmv{} is an extension of \nusmv{} and is capable of checking synchronous finite and infinite state systems \cite{DBLP:conf/cav/CavadaCDGMMMRT14}.
\nuxmv{} has its own input format based on the language for \nusmv{}, but it is also possible to pass directly \vmtlib{} thanks to its \texttt{read\_vmt\_model} command.
Several algorithms are available for model checking and both invariants and live properties are supported.
The logic supported by \nuxmv{} depends on the chosen algorithm since algorithms working with interpolants do not support having both integer and real arithmetic at the same time.
Complex logics containing arrays, BitVectors, integer and real arithmetic are otherwise supported.

Options are passed to the solver through the standard input as commands, the output is parsed through the use of regular expressions.

\paragraph*{\iceia{}.}
\iceia{} utilizes implicit abstraction to verify the safety of properties over infinite state systems \cite{10.1007/978-3-319-41528-4_15}.
It supports \vmtlib{} directly since it's the only format which it reads, and allows for both invariant and live properties.
\iceia{} supports linear arithmetic formulae which can contain integer and real arithmetic combined through boolean conditions.

The output from the solver may contain a counterexample encoded as a series of steps, each of which is a conjunction of assignments in \smtlib{} format. To parse the counterexample, the \texttt{SmtLibParser} class is employed.

\paragraph*{\euforia{}.}
\euforia{} is a model checker which works through abstraction-refinement to prove a property or find a feasible abstract counterexample that disproves it \cite{10.1007/978-3-030-11245-5_17}.
The \vmtlib{} format is supported directly by \euforia{} as its input format and it allows for invariant properties.
\euforia{} supports formulae containing BitVectors and arrays, it does not on the other hand support linear or real arithmetic.

For unsafe properties a counterexample is given. It consists of a list of function declarations each corresponding to a symbol in the original model and the index of a step, they contain the assignment to the symbol in the specified step. Again, the \texttt{SmtLibParser} is used to parse the result.


\section{Results}
Results from model checking are parsed from \pyvmt{} and returned as an object of class \texttt{Result}.
This allows accessing the information regarding the produced output without requiring parsing it on the user's side.

Typically results specify whether the checked property is safe or unsafe, in the case of an unsafe property a counterexample in the form of a trace may be produced.

\section{Traces}
Whenever a property is not verified, a model checker may provide a counterexample corresponding to a series of steps that lead the system to an unsafe state.
The output from different model checkers is usually not consistent, but \pyvmt{} offers a uniform interface for reading a trace produced by any of the supported model checkers.

In \pyvmt{} traces are modelled as a series of steps, each of which consists of a mapping between the state variables and their values at the step.
Traces may also contain an infinite amount of steps (for example when working with live properties).
One of the possible modelling approaches is used in which exactly one of the steps is a loopback step: after the last step in the series the trace loops back to this step and continues infinitely.

\paragraph*{Iterating through a trace.}
\pyvmt{} offers many methods for reading the information relating to a trace, it's possible to retrieve every step using the \texttt{get\_steps} method, or a specific step using \texttt{get\_step} method passing the step index.

From each step, the previous and the next step are also available using the \texttt{get\_next\_step} and \texttt{get\_prev\_step} methods.

\paragraph*{Comparing steps.}
It's not always important to read every variable on each step, sometimes just the ones which are changing or which have changed since the last step may be more useful.
The \texttt{TraceStep} class implements methods to perform this comparisons automatically which are \texttt{get\_changing\_variables} and \texttt{get\_changed\_variables}.

Comparing two specific steps is also possible thanks to the \texttt{get\_different\_variables} method of the Step.

\paragraph*{Loopback step.}
For traces with infinite steps a loopback step is present. A trace can only have one loopback step and it can be obtained with the \texttt{get\_loopback\_step} method.

\begin{listing}
    \label{alg:model-checking-traces}
    \inputmintedpy{py/model_checking.py}{22}{46}
    \caption{The trace from the previous example is printed and analyzed in various ways, \pyvmt{} allows for a variety of operations to access the trace and compare the steps.}
\end{listing}

\paragraph*{Evaluating formulae.}
To analyze a counterexample it may be useful to evaluate some formula over certain steps of the trace, or over pairs of states to verify the behavior during transitions.

To aid this kind of computation the \texttt{TraceStep} class contains an \texttt{evaluate\_formula} method which takes as input some arbitrary \pysmt{} formula and evaluates it.
To use this feature over a transition it is also possible to pass a formula containing the \texttt{Next} operator, the value for a next state variable will then be taken from the next step.

The evaluation consists of the following steps:
\begin{enumerate}
    \item \texttt{Next} operators are pushed to the leaves to allow for substitution of next state variables
    \item Substitutions for the symbols are produced, from current and next step if available
    \item The calculated substitutions are applied over the formula
    \item The \texttt{FormulaSimplifier} is used in order to evaluate the formula
\end{enumerate}

If the formula can be simplified completely the result will simply be a single \texttt{FNode} containing a constant value.

\begin{listing}[H]
    \label{alg:model-checking-evaluation}
    \inputmintedpy{py/formula_evaluation.py}{30}{53}
    \caption{Model containing two counters, one from 0 to 3 and one from 0 to 2. The property $\G (\mathrm{a} + \mathrm{b} < 5)$ is unsafe, and formula evaluation is used to analyze the resulting trace.}
\end{listing}


      \chapter{\ltl{} Support}
\label{ch:ltl-support}
\ltl{} properties aren't always supported directly by model checkers, since usually just safety and liveness properties are supported, which are also included within the core language of \vmtlib{}.
The importance of these two kinds of properties cannot be overstated, in fact their wide adoption is justified by the fact that they're sufficient even to verify more complex properties like those that can be written in \LTL{} as long as they're encoded properly.

Over the years multiple algorithms were developed to tackle this problem, two of them were chosen to be implemented as part of the thesis work.
These functionalities aim to assist the end user, offering a layer of abstraction that simplifies the model checking workflow, but they also serve as proof of concept to display the capabilities of this library and how it may be used to create reusable tools and implement complex algorithms in \python{} that don't depend on any underlying solver thanks to \pysmt{}'s formula manipulation tools.

\section{\ltltosmv{} Encoding}
One of the classic encoding procedures for \ltl{} was proposed in 1994 in a paper titled "Another Look at \ltl{} Model Checking" \cite{DBLP:conf/cav/ClarkeGH94}.
This procedure involves creating a tableau associated with the negation of the \ltl{} formula, finding a set of elementary subformulas and for each of them a condition that defines whether that subformula is satisfied (characteristic function, here referred to as sat value).
By composing the model with the tableau it's possible to analyze how the state of the \ltl{} specification changes as steps are performed on the original model.

Elementary subformulae are defined for the \ltl{} operators \fX{} and \fU{}, therefore the first step for encoding is normalizing the negated formula through the \texttt{LtlRewriter} walker which performs the following substitutions:
\begin{itemize}
    \item \fR{} $\longrightarrow{}$ $\lnot(\lnot \phi \U \lnot \psi)$
    \item \fF{} $\longrightarrow{}$ $\top \U \phi$
    \item \fG{} $\longrightarrow{}$ $\lnot (\F \lnot \phi)$
\end{itemize}

A specialized formula walker, \texttt{LtlEncodingWalker}, is then employed to compute all of the elementary subformulae and the associated sat values:
\begin{itemize}
    \item For \fX{} a fresh symbol \elx{} is created, with $\mathrm{sat}(\fX) = \elx{}$
    \item For \fU{} a fresh symbol \elu{} is created, with $\mathrm{sat}(\fU) = sat(\psi) \lor (\mathrm{sat}(\phi) \land \elu)$
\end{itemize}

To constrain the evolution of the new variables, a transition constraint is added for each of them:

\[ \elx{} \leftrightarrow (\mathrm{sat}(\phi))' \]

To ensure that the properties are eventually fulfilled, liveness properties are added to the model.
For every elementary subformula in the shape $\X (\fU)$, a property $\lnot \mathrm{sat}(\fU) \lor \mathrm{sat}(\psi)$ is added.
The properties are then combined into a single liveness property to ensure compatibility with the model checkers, which is negated to return to our original \ltl{} specification.

% Include sample implementation?

\section{\ltl{} Circuit Encoding}

A different approach for encoding was presented in "A circuit Approach to \ltl{} Model Checking" \cite{DBLP:conf/fmcad/ClaessenES13}.
This encoding works a little differently since it doesn't generate the tableau corresponding to the negated \ltl{} formula, it instead creates circuit-like monitors for each of the subformulae in the negated \ltl{} formula capable of verifying whether a counterexample to the specification exists.

For example the specification $\X a \land (b \lor \F \lnot c) \U d$ is negated into $\X \lnot a \lor (\lnot b \land \G c) \R \lnot d$.
The formula is split into its subformulae, each of which is denoted by a variable creating an equivalent formula:
\begin{align}
           &     \mathit{z_0} \\
    \land\ & \G (\mathit{z_0} \rightarrow G c) \\
    \land\ & \G (\mathit{z_1} \rightarrow \lnot b \land \mathit{z_0}) \\
    \land\ & \G (\mathit{z_2} \rightarrow \mathit{z_1} \R \lnot d) \\
    \land\ & \G (\mathit{z_3} \rightarrow \X \lnot a) \\
    \land\ & \G (\mathit{z_4} \rightarrow \mathit{z_3} \lor \mathit{z_2})
\end{align}

For the process to work, the negated formula has to be yet again normalized, this time in \nnf{} or \NNF{}. \pysmt{} already has a formula walker to perform this rewriting but it doesn't support the custom \ltl{} operators, therefore an extension was required to add the following rewritings recursively:
\begin{itemize}
    \item $\lnot \fX$ $\longrightarrow$ $\X \lnot \phi$
    \item $\lnot \fG$ $\longrightarrow$ $\F \lnot \phi$
    \item $\lnot \fF$ $\longrightarrow$ $\G \lnot \phi$
    \item $\lnot (\fU)$ $\longrightarrow$ $\lnot \phi \R \lnot \psi$
    \item $\lnot (\fR)$ $\longrightarrow$ $\lnot \phi \U \lnot \psi$
\end{itemize}

A specialized formula walker called \texttt{LtlCircuitEncodingWalker} is used to find all of the elementary subformulae.
For every \ltl{} operator, as well as the \texttt{And} and \texttt{Or} operators, a \texttt{FreshSymbol} is created: this is the activator for the monitor, once it's verified the monitor will start monitoring the formula and outputting its state.
The output for this walker is the list of activators and the formula that the monitor observes.

Monitors have 3 outputs: \texttt{accept}, \texttt{failed}, and \texttt{pending}.
\begin{itemize}
    \item \texttt{pending} is used by the monitor to specify that it's waiting for a signal from the inputs;
    \item \texttt{failed} indicates that a violation of the specification has been detected, meaning that the trace is no longer a valid counterexample;
    \item If \texttt{accept} goes false forever it's an indication that the trace is not a valid counterexample.
\end{itemize}

The \texttt{LtlCircuitEncodingWalker} has a \texttt{make\_monitor} method that generates the outputs for the monitor, as well as the state variables and their constraints required by the monitor to function.
The exact values for the signals may be found in the mentioned paper, below the section `List of Monitors for \pltl'.
The initial activator $\mathit{z_0}$ is replaced by the variable $\mathit{is\_init}$ which is modelled to be true in the first cycle, thus the circuit starts observing the state of the specification right at the start of the model checking process.

After all of the monitors have been created there are just two more things to compute, the \texttt{has\_failed} condition which is set to true indefinitely when any of the monitors fails, and the \texttt{accept} condition which specifies that every monitor must output \texttt{accept} infinitely often in order for the trace to be a valid counterexample.
Finally, a single liveness property is generated, $\lnot \mathit{accept}$.
If this property is demonstrated to be true it means that a counterexample doesn't exist, which proves our original specification.

% Include sample implementation?

\section{Experimental Evaluation}

To validate the efficacy of this approach, an experimental evaluation was conducted on a set of benchmarks.
The time required by the encoding procedure is negligible compared to the time required by a solver to verify the validity of the generated property, what this evaluation focussed on is therefore the quality of the generated encoding in terms of number of extra state variables required, and the time for solving.

\begin{figure}
    \includegraphics{ltlenc_time_comparison.png}
    \caption{Comparison between LTL2SMV and Circuit encoding in terms of time required for model checking.}
    \label{fig:ltlenc_time_comparison}
\end{figure}

The encoders were ran on a set of almost 600 benchmarks, after the conversion \iceia{} was run to evaluate the efficacy of the encoding in terms of aiding the model checking process.
From the output of the solver, the \texttt{total\_time} was collected and used as a metric for comparison.
On the set of benchmarks which was analyzed, the \ltltosmv{} encoding generally performed better than the circuit encoding version, this can be seen in the Figure \ref{fig:ltlenc_time_comparison}.

\begin{figure}
    \includegraphics{ltlenc_stvar_comparison.png}
    \caption{Comparison between LTL2SMV and Circuit encoding in terms of number of supplementary variables required to encode the property.}
    \label{fig:ltlenc_stvar_comparison}
\end{figure}

A possible explanation as to why the circuit encoding is slower can be found in Figure \ref{fig:ltlenc_stvar_comparison} which shows that this kind of encoding with the current implementation requires a much larger number of state variables compared to the \ltltosmv{} version.
It's possible to improve the efficiency of this algorithm by drastically reducing the number of auxiliary state variables required as suggested in its paper, by not creating a new state variable for each of the \texttt{And} and \texttt{Or} logic operators, instead using a single activator for fomulae containing multiple of them.


      \chapter{Other Features}
\label{ch:other-features}

This chapter will discuss in detail some of the other features that are available within \pyvmt{}, particularly focussing on tools which may be useful or even necessary for the end user to perform certain computations while using a model.

\section{Merging Transition Systems}
\paragraph*{Synchronous compositions.}
Composing transition systems by means of synchronous composition involves having the two systems evolve through transitions at the same time.
This differs from asynchronous composition where at each step, exactly one of the systems is chosen to evolve independently of the other.
For symbolic transition systems, synchronous composition means creating a transition system that includes all of the state and input variables of the original systems, and constraining the system so that the initial states are valid initial states for both transition systems, and the transitions are valid transitions for both transition systems.
This can simply be done by means of conjunction between the original constraints.

Merging symbolic transition systems through synchronous composition in \pyvmt{} can be done using the \texttt{compose} function, which takes as input two models and returns the resulting model.
It was decided to keep the functionality as simple as possible, in order to avoid behaviors which the user might not expect like having the variables change to \texttt{FreshSymbol}s as the systems are merged.
For that reason, whenever a state variable and/or an input coincide, they're considered the same variable in the final system.
This can be useful to make the two models communicate by replacing a nondeterministic input variable with a constrained state variable as defined from another model.

\begin{listing}
    \label{alg:composition}
    \inputmintedpy{py/composition.py}{9}{33}
    \caption{\texttt{model\_b} is a counter which increments based on a non-deterministic variable \texttt{a}. Since \texttt{a} may always be 0, the property $\F (\G (\mathrm{counter} > 10))$ is unsafe.
    Composing the model with \texttt{model\_a} replaces the input \texttt{a} with the state variable from \texttt{model\_a}, making the previous property safe.}
\end{listing}

\section{Renaming Transition Systems}
It may be necessary to rename the variable contained within a model to avoid having an overlap with the variables of another model when using synchronous composition.
\pyvmt{} offers a model renamer to simplify such a process, with several utility functions that implement a specific renaming pattern.
It's also possible for the user to specify a custom renaming callback by using the \texttt{rename} function.
The patterns which are available act on the prefix and suffix of each of the input and state variables, allowing the addition or replacement of the prefix/suffix.


\begin{listing}[H]
    \label{alg:renaming}
    \inputmintedpy{py/renaming.py}{9}{41}
    \caption{Through renaming, 3 counter models are created, each with a different limit.
    The property $\G (\sum_\textrm{i}^3 \mathrm{model[i].a} < 6)$ is then checked and a counterexample of length 11 is found, when the counters synchronize and their sum reaches 6.}
\end{listing}


      \chapter{Conclusion}
\label{cha:conclusions}



      %\input{capitolo4}
      
      
    \endgroup


    % bibliografia in formato bibtex
    %
    % aggiunta del capitolo nell'indice
    \addcontentsline{toc}{chapter}{Bibliography}
    % stile con ordinamento alfabetico in funzione degli autori
    \bibliographystyle{plain}
    \bibliography{biblio}
%%%%%%%%%%%%%%%%%%%%%%%%%%%%%%%%%%%%%%%%%%%%%%%%%%%%%%%%%%%%%%%%%%%%%%%%%%
%%%%%%%%%%%%%%%%%%%%%%%%%%%%%%%%%%%%%%%%%%%%%%%%%%%%%%%%%%%%%%%%%%%%%%%%%%
%% Nota
%%%%%%%%%%%%%%%%%%%%%%%%%%%%%%%%%%%%%%%%%%%%%%%%%%%%%%%%%%%%%%%%%%%%%%%%%%
%% Nella bibliografia devono essere riportati tutte le fonti consultate 
%% per lo svolgimento della tesi. La bibliografia deve essere redatta 
%% in ordine alfabetico sul cognome del primo autore. 
%% 
%% La forma della citazione bibliografica va inserita secondo la fonte utilizzata:
%% 
%% LIBRI
%% Cognome e iniziale del nome autore/autori, la data di edizione, titolo, casa editrice, eventuale numero dell’edizione. 
%% 
%% ARTICOLI DI RIVISTA
%% Cognome e iniziale del nome autore/autori, titolo articolo, titolo rivista, volume, numero, numero di pagine.
%% 
%% ARTICOLI DI CONFERENZA
%% Cognome e iniziale del nome autore/autori (anno), titolo articolo, titolo conferenza, luogo della conferenza (città e paese), date della conferenza, numero di pagine. 
%% 
%% SITOGRAFIA
%% La sitografia contiene un elenco di indirizzi Web consultati e disposti in ordine alfabetico. 
%% E’ necessario:
%%   Copiare la URL (l’indirizzo web) specifica della pagina consultata
%%   Se disponibile, indicare il cognome e nome dell’autore, il titolo ed eventuale sottotitolo del testo
%%   Se disponibile, inserire la data di ultima consultazione della risorsa (gg/mm/aaaa).    
%%%%%%%%%%%%%%%%%%%%%%%%%%%%%%%%%%%%%%%%%%%%%%%%%%%%%%%%%%%%%%%%%%%%%%%%%%
%%%%%%%%%%%%%%%%%%%%%%%%%%%%%%%%%%%%%%%%%%%%%%%%%%%%%%%%%%%%%%%%%%%%%%%%%%
    

    % \titleformat{\chapter}
    %     {\normalfont\Huge\bfseries}{Allegato \thechapter}{1em}{}
    % sezione Allegati - opzionale
    % \appendix
    % \input{allegati}

\end{document}
